% GENERATED FILE. Edit canonical lessons, then run npm run generate:books.
% canonical-source-hash: fnv1a64:d716507f185675b2
% canonical-lessons: FR-C206-casser, FR-C206-toucher, FR-C206-gouter, FR-C206-chauffer, FR-C206-refroidir

\chapter{To break, To touch, To taste, To heat, To cool down}
\label{ch:fr-to-break-to-touch-to-taste-to-heat-to-cool-down}
\hlchaptermodality{206}

\begin{chapteropening}
\textbf{By the end of this chapter:} \emph{I can say to break, to touch, to taste, to heat and to cool down.}

Complete the last lesson of chapter 206: I can say to break, to touch, to taste, to heat and to cool down.
\end{chapteropening}

\section[casser]{casser — to break}

\label{lesson:FR-C206-casser}



\begin{quote}
Before the new one: say the French for to order, then the French for to dry.
\end{quote}

\subsection*{You'll want to know: casser}
\textbf{casser} — \textquotedblleft{}to break\textquotedblright{}.

\begin{cousinweb}
\textbf{casser} $\leftarrow$ Old French \emph{quasser} $\leftarrow$ Latin \emph{quassare}, \textquotedblleft{}to shake hard, to shatter\textquotedblright{} — inherited. (EtymDB 2.1)
\end{cousinweb}

\begin{grammarlens}[title={What you've built}]
One more word for this chapter.
\end{grammarlens}

\subsection*{Your turn}
\begin{itemize}
\raggedright
  \item \emph{Say it:} \emph{casser}
  \item \emph{Say it:} \emph{casser}, once more
  \item \emph{Say it:} say \emph{casser}
  \item \emph{Recall:} say the French for to order, then the French for to dry, then say \emph{casser} again
\end{itemize}

\begin{tcolorbox}[breakable,colback=teal!4,colframe=teal!35!black,title={Before you move on}]
What does casser mean? (\textquotedblleft{}To break\textquotedblright{}.) Say it once more.
\end{tcolorbox}
\section[toucher]{toucher — to touch}

\label{lesson:FR-C206-toucher}



\begin{quote}
Before the new one: say the French for to dry, then the French for to break.
\end{quote}

\subsection*{You'll want to know: toucher}
\textbf{toucher} — \textquotedblleft{}to touch\textquotedblright{}.

\begin{cousinweb}
\textbf{toucher} $\leftarrow$ Vulgar Latin \emph{*toccare}, \textquotedblleft{}to knock\textquotedblright{}, imitating the sound of a tap — inherited; the account is the usual one, not certain.
\end{cousinweb}

\begin{grammarlens}[title={What you've built}]
One more word for this chapter.
\end{grammarlens}

\subsection*{Your turn}
\begin{itemize}
\raggedright
  \item \emph{Say it:} \emph{toucher}
  \item \emph{Say it:} \emph{toucher}, once more
  \item \emph{Say it:} say \emph{toucher}
  \item \emph{Recall:} say the French for to dry, then the French for to break, then say \emph{toucher} again
\end{itemize}

\begin{tcolorbox}[breakable,colback=teal!4,colframe=teal!35!black,title={Before you move on}]
What does toucher mean? (\textquotedblleft{}To touch\textquotedblright{}.) Say it once more.
\end{tcolorbox}
\section[goûter]{goûter — to taste}

\label{lesson:FR-C206-gouter}



\begin{quote}
Before the new one: say the French for to break, then the French for to touch.
\end{quote}

\subsection*{You'll want to know: goûter}
\textbf{goûter} — \textquotedblleft{}to taste\textquotedblright{}.

\begin{cousinweb}
\textbf{goûter} $\leftarrow$ Old French \emph{goster} $\leftarrow$ Latin \emph{gustare}, \textquotedblleft{}to taste\textquotedblright{} (EtymDB 2.1) — inherited; the circumflex marks the lost \emph{s}.
\end{cousinweb}

\begin{grammarlens}[title={What you've built}]
One more word for this chapter.
\end{grammarlens}

\subsection*{Your turn}
\begin{itemize}
\raggedright
  \item \emph{Say it:} \emph{goûter}
  \item \emph{Say it:} \emph{goûter}, once more
  \item \emph{Say it:} say \emph{goûter}
  \item \emph{Recall:} say the French for to break, then the French for to touch, then say \emph{goûter} again
\end{itemize}

\begin{tcolorbox}[breakable,colback=teal!4,colframe=teal!35!black,title={Before you move on}]
What does goûter mean? (\textquotedblleft{}To taste\textquotedblright{}.) Say it once more.
\end{tcolorbox}
\section[chauffer]{chauffer — to heat}

\label{lesson:FR-C206-chauffer}



\begin{quote}
Before the new one: say the French for to touch, then the French for to taste.
\end{quote}

\subsection*{You'll want to know: chauffer}
\textbf{chauffer} — \textquotedblleft{}to heat\textquotedblright{}.

\begin{cousinweb}
\textbf{chauffer} $\leftarrow$ Old French \emph{chauffer} (EtymDB 2.1) $\leftarrow$ Vulgar Latin \emph{*calefare}, from Latin \emph{calefacere}, \textquotedblleft{}to make warm\textquotedblright{} (Lewis \& Short) — inherited.
\end{cousinweb}

\begin{grammarlens}[title={What you've built}]
One more word for this chapter.
\end{grammarlens}

\subsection*{Your turn}
\begin{itemize}
\raggedright
  \item \emph{Say it:} \emph{chauffer}
  \item \emph{Say it:} \emph{chauffer}, once more
  \item \emph{Say it:} say \emph{chauffer}
  \item \emph{Recall:} say the French for to touch, then the French for to taste, then say \emph{chauffer} again
\end{itemize}

\begin{tcolorbox}[breakable,colback=teal!4,colframe=teal!35!black,title={Before you move on}]
What does chauffer mean? (\textquotedblleft{}To heat\textquotedblright{}.) Say it once more.
\end{tcolorbox}
\section[refroidir]{refroidir — to cool down}

\label{lesson:FR-C206-refroidir}



\begin{quote}
Before the new one: say the French for to taste, then the French for to heat.
\end{quote}

\subsection*{You'll want to know: refroidir}
\textbf{refroidir} — \textquotedblleft{}to cool down\textquotedblright{}.

\begin{cousinweb}
\textbf{refroidir} — built from \emph{re-} and the word for \textquotedblleft{}cold\textquotedblright{}, $\leftarrow$ Latin \emph{frigidus} (EtymDB 2.1), with the verb ending \emph{-ir}.
\end{cousinweb}

\begin{grammarlens}[title={What you've built}]
One more word for this chapter.
\end{grammarlens}

\subsection*{Your turn}
\begin{itemize}
\raggedright
  \item \emph{Say it:} \emph{refroidir}
  \item \emph{Say it:} \emph{refroidir}, once more
  \item \emph{Say it:} say \emph{refroidir}
  \item \emph{Recall:} say the French for to taste, then the French for to heat, then say \emph{refroidir} again
\end{itemize}

\begin{tcolorbox}[breakable,colback=teal!4,colframe=teal!35!black,title={Before you move on}]
What does refroidir mean? (\textquotedblleft{}To cool down\textquotedblright{}.) Say it once more.
\end{tcolorbox}
